\begin{claim}{D.35}{Proved}{One exact conjugator equation and its finite obstructions}
Let $P$ be the full pure three-cusp group of D.32, let
$F=\widehat F(x,y)$, and put
$\delta(x)=x^{-1}$, $\delta(y)=y^{-1}$.
For a nonidentity involution $\alpha\in P$, use its normalized
lift from D.34,
\[
 \alpha(x)=x^{-1},\qquad
 \alpha(y)=f^{-1}y^{-1}f,\qquad f\in F'.
\]
Define the compact solution space
\[
\begin{split}
 \mathcal S_\alpha=\{(\mu,h,u)\in
       \widehat\Z^\times\times F'\times F:\quad&
       h^{-1}\alpha(h)=f,\\
       &x^\mu h y^\mu h^{-1}
                    =u(xy)^\mu u^{-1}\}.
\end{split}                                                   \tag{D.35a}
\]
The assignment
\[
 (\mu,h,u)\longmapsto
 \left[\beta_{\mu,h}\right],\qquad
 \beta_{\mu,h}(x)=x^\mu,\quad
 \beta_{\mu,h}(y)=h y^\mu h^{-1},                              \tag{D.35b}
\]
has image exactly
\[
 \operatorname{Conj}_{P}(\delta,\alpha)
       =\{\beta\in P:\ \beta\delta\beta^{-1}=\alpha\}.
                                                                  \tag{D.35c}
\]
Every map (D.35b) is an automorphism of the full profinite
group.  The fibers of the assignment are exactly
$(\mu,h,u)\mapsto(\mu,h,u(xy)^a)$ for $a\in\widehat\Z$.
It therefore induces a homeomorphism from the corresponding
quotient of $\mathcal S_\alpha$ to (D.35c).

For $\alpha\in GT_0$, the inclusion $GT_0\hookrightarrow P$
gives the equality of actual conjugator sets
\[
       \operatorname{Conj}_{GT_0}(\delta,\alpha)
             =\operatorname{Conj}_{P}(\delta,\alpha).             \tag{D.35d}
\]
The defining Harbater--Schneps relations for the conjugating
pair follow from the proof below; they are not discarded.
Finally, the existence of a solution of (D.35a) is equivalent
to the nonemptiness of its explicit finite solution sets
(D.35h) at every characteristic finite quotient of $F$.
\end{claim}
\begin{proof}
For $(\mu,h,u)\in\mathcal S_\alpha$, the free universal
property first gives a continuous endomorphism
$\beta_{\mu,h}:F\to F$ with the displayed generator images.
Write $H=\beta_{\mu,h}(F)$, a closed subgroup because $F$
is compact.  It is generated by two elements, so
$\dim H^1(H,\F_2)\leq2$.  The restrictions of the two
standard characters of $F$ to these generators are the
two coordinate characters: $\mu\bmod2=1$ and conjugation
does not change abelianization.  These restrictions are
independent, proving
\[
                         \dim H^1(H,\F_2)=2.                   \tag{D.35e}
\]
The unit exponent $\mu$ shows that $H$ contains
$\overline{\langle x\rangle}$ and
$h\overline{\langle y\rangle}h^{-1}$.
The second equation in (D.35a) puts
$u\overline{\langle xy\rangle}u^{-1}$ in $H$ as well.
In particular $H$ contains the three full cusp
$\Z_2$ subgroups required by D.32.  Its closed-subgroup
rank test, applied to (D.35e), gives $H=F$.
A surjective endomorphism of a finitely generated
profinite group is injective, so $\beta_{\mu,h}$ is
an automorphism.  Its images of the three cusps,
including the exact displayed equation for the third,
place its outer class in $P$.

The first equation in (D.35a) is $\alpha(h)=hf$.
Substitution on the two free generators gives
\[
 \alpha\beta_{\mu,h}(x)=x^{-\mu}
                          =\beta_{\mu,h}\delta(x),\qquad
 \alpha\beta_{\mu,h}(y)=h y^{-\mu}h^{-1}
                          =\beta_{\mu,h}\delta(y).
                                                                  \tag{D.35f}
\]
These are equalities of actual automorphisms, proving
that the image of (D.35b) lies in (D.35c).

Conversely, let $\beta\in P$ satisfy (D.35c), and take its
unique normalized representative from D.32.
Write it as $\beta_{\mu,h}$ with $h\in F'$.
The normalization is a homomorphic section
$P\to\Aut(F)$, so the outer conjugacy equality lifts
to the actual equality $\alpha\beta_{\mu,h}
=\beta_{\mu,h}\delta$.
Its value on $y$ implies that
$h^{-1}\alpha(h)f^{-1}$ centralizes $y^{-\mu}$.
By A.7 it is a power of $y$.  It also lies in $F'$,
so its abelianization forces this power to be zero.
Thus $h^{-1}\alpha(h)=f$.
The third-cusp condition on $\beta$ gives a conjugator
$u$ in the second equation of (D.35a).
This proves that the image is exactly (D.35c).

The normalized pair $(\mu,h^{-1})$ is unique.
For this fixed pair the two third-cusp conjugators
$u,u'$ differ by right multiplication by the
centralizer of $(xy)^\mu$, which is
$\overline{\langle xy\rangle}$ by A.7.
These are precisely the asserted fibers.
The assignment is continuous in the profinite
automorphism topology, and its domain is compact.
Its quotient is therefore homeomorphic to its
Hausdorff image.

For an alternative form using D.34, fix
$h_0\in F'$ with $h_0^{-1}\alpha(h_0)=f$, and put
$C=C_F(\alpha)$.
All solutions of the first equation of (D.35a)
are exactly $h=ch_0$ with $c\in C$:
$\alpha(h)=hf$ and $\alpha(h_0)=h_0f$ give
$\alpha(hh_0^{-1})=hh_0^{-1}$, and the reverse
substitution proves the converse.
Since $\alpha$ acts by $-1$ on $F^{\rm ab}$,
every $c\in C$ has $2[c]=0$, hence $[c]=0$.
Thus $C\subset F'$ and this substitution respects
the stated domain of (D.35a).  The remaining exact
equation is
\[
 x^\mu c h_0 y^\mu h_0^{-1}c^{-1}
                         =u(xy)^\mu u^{-1}.
                                                                  \tag{D.35g}
\]

We prove (D.35d) by identifying the relevant centralizer,
rather than assuming that a conjugator in $P$ lies in
$GT_0$.  Let $P^1=\ker(\lambda:P\to\widehat\Z^\times)$.
First,
$C_{P^1}(\delta)=1$.
Indeed a normalized element $(1,k)\in P^1$ commuting
with $\delta$ has $k=\delta(k)$, by uniqueness of the
normalized representative and its group law.
The following explicit mutually inverse homomorphisms
identify the semidirect product:
\[
 \begin{aligned}
 C_2\amalg C_2\amalg C_2
    &\xrightarrow{\ \sim\ }F\rtimes_\delta\langle s\rangle,
       & (s_0,s_1,s_2)&\longmapsto(s,xs,ys),\\
 (x,y,s)&\longmapsto(s_1s_0,s_2s_0,s_0).
 \end{aligned}
\]
The orders and conjugation relations follow by direct
substitution.  A free profinite factor has trivial
intersection with a conjugate by an element outside
that factor; the original full-profinite statement is
recalled in \cite[Theorem 5 of the inspected author
version]{A-HerfortZalesskii2008}.
Consequently the centralizer of $s_0$ in this free
product is $\langle s_0\rangle$.
Its intersection with $F$ is trivial.
Hence $C_F(\delta)=1$, so $k=1$ and
$C_{P^1}(\delta)=1$ as claimed.

Let $\Sigma\cong S_3$ be the geometric cusp-permutation
subgroup of the full loop-outer group, generated by
the outer classes
$\theta(x)=y,\ \theta(y)=x$ and
$\omega(x)=y,\ \omega(y)=(xy)^{-1}$.
The exact identification
$GT_0=C_P(\Sigma)$ follows from
\cite[Propositions 3--4 and Theorem 2]{R-HS}:
commutation with $\theta$ is relation (I), and,
after (I), commutation with $\omega$ is relation (II).
Both $\alpha$ and $\delta$ centralize $\Sigma$.
For any $\beta\in\operatorname{Conj}_P(\delta,\alpha)$
and $\sigma\in\Sigma$, form the actual element of
the outer group
$r_\sigma=\beta^{-1}\sigma\beta\sigma^{-1}$.
The pure subgroup $P$ is normal in the full
loop-outer group, and the common multiplier is a
homomorphism there.  Therefore $r_\sigma\in P^1$.
Moreover
\[
 \begin{split}
 r_\sigma\delta r_\sigma^{-1}
  &=\beta^{-1}\sigma\beta
        (\sigma^{-1}\delta\sigma)
        \beta^{-1}\sigma^{-1}\beta\\
  &=\beta^{-1}\sigma\alpha\sigma^{-1}\beta
    =\beta^{-1}\alpha\beta=\delta.
 \end{split}
\]
The proved equality $C_{P^1}(\delta)=1$ gives
$r_\sigma=1$.  Thus $\beta$ centralizes $\Sigma$,
so its normalized pair satisfies both (I) and (II)
and lies in $GT_0$.  This proves (D.35d).

Here is the exact finite form of the remaining
existence problem.  For each open characteristic
subgroup $N\triangleleft F$, put $Q_N=F/N$ and
$m_N=\exp(Q_N)$.
The automorphism $\alpha$ induces $\alpha_N$.
Use bars for the images of $x,y,f$ in $Q_N$ and define
\[
\begin{split}
 \mathcal S_{\alpha,N}=\{(a,b,v)\in
       (\Z/m_N)^\times\times Q_N'\times Q_N:\quad&
       b^{-1}\alpha_N(b)=\bar f,\\
       &\bar x^a b\bar y^a b^{-1}
              =v(\bar x\bar y)^a v^{-1}\}.
\end{split}                                                       \tag{D.35h}
\]
These are finite sets, obtained entirely from the
specified finite quotient and its induced
automorphism.  Reduction sends
$\mathcal S_\alpha$ to every
$\mathcal S_{\alpha,N}$.
Conversely, a point of a nonempty
$\mathcal S_{\alpha,N}$ lifts to
$\widehat\Z^\times\times F'\times F$:
units lift by the Chinese remainder theorem,
and the image of $F'$ in $Q_N$ is exactly $Q_N'$.
The lifts satisfying the two equations modulo $N$
form a nonempty closed subset of this compact space.
Every finite collection of such conditions is
implied by the one at a common characteristic
refinement.  Their finite intersection property
therefore gives a point satisfying both equations
in $F$ itself.  Open characteristic subgroups are
cofinal for a finitely generated profinite group,
so equality at all these quotients is equality in
$F$.  This proves the asserted finite criterion.

The need for the full third-cusp equation can be
checked by a short executed finite certificate.
Put $d=[x,y]=xyx^{-1}y^{-1}$ and take the actual
commutator word
\[
                 w=y d^{-1}y^{-1}x^{-1}d^{-1}x\in F'.
\]
Consider the endomorphism
$\rho(x)=x$, $\rho(y)=wyw^{-1}$.
Its abelianization is the identity, so the Frattini
argument makes its induced endomorphism on every
maximal pro-$p$ quotient an automorphism.
In the finite quotient
\[
 F\twoheadrightarrow A_5,\qquad
             x\longmapsto(1\,2\,3),\quad
             y\longmapsto(3\,4\,5),
\]
the program
\path{code/part_a/pure_involution_conjugator.py}
enumerates the generated group of order $60$ and
computes
\[
 \bar w=(1\,4\,3\,2\,5),\qquad
 \overline{wyw^{-1}}=(1\,2\,3).
\]
The image of $\rho(F)$ in this quotient therefore
has order $3$, proving $\rho(F)\ne F$.
Moreover $\overline{xy}$ has order $5$, whereas
$\overline{xwyw^{-1}}$ has order $3$; the
third-cusp equation fails in this same quotient.
This word is a certificate concerning the
generation test.  It is not asserted to solve
the coboundary equation for a nonstandard
involution.  No nonemptiness or emptiness
assertion for the spaces (D.35a) of arbitrary
involutions is inferred from this finite example.
\end{proof}
